\section{2022级工科数学分析下 B卷}

\subsection{试题}

\paragraph*{一、填空题：共 5 题，每题 2 分，共 10 分。}
\begin{enumerate}
  \item 微分方程 $y''+3y'+2y=0$ 的通解为\underline{\hspace{4cm}}。
  \item 函数 $u=2xy-z^2$ 在点 $(2,-1,-1)$ 处方向导数的最大值为\underline{\hspace{4cm}}。
  \item 设 $\Gamma$ 是圆周 $x^2+y^2=1$ 在第一象限的部分，则第一类曲线积分
  \[
    \int_\Gamma (x+y)^2\,ds=\underline{\hspace{4cm}}.
  \]
  \item 级数 $\displaystyle \sum_{n=0}^{\infty}(-1)^n\dfrac{2n+3}{(2n+1)!}$ 的和为\underline{\hspace{4cm}}。
  \item 设周期为 $2\pi$ 的函数
  \[
    f(x)=
    \begin{cases}
      x^2, & -\pi<x\le 0,\\
      0, & 0<x\le \pi,
    \end{cases}
  \]
  则 $f(x)$ 的傅里叶（Fourier）级数在 $x=-\pi$ 处收敛于\underline{\hspace{4cm}}。
\end{enumerate}

\paragraph*{二、选择题：共 5 题，每题 2 分，共 10 分。}
\begin{enumerate}
  \item 下列微分方程中，属于二阶线性常微分方程的是（\quad）
  \begin{enumerate}
    \item[A.] $y''+x^3y'+(\ln x)y=\tan x$；
    \item[B.] $(x^2+y^2)\,dy+y^2\,dx=0$；
    \item[C.] $(y'')^2+x^2y=1$；
    \item[D.] $y''+2\ln y=2x$。
  \end{enumerate}
  \item 已知函数 $f(x)$ 具有二阶连续导函数，且 $f(x)>0,\ f'(0)=0$（$f$ 在 $0$ 处的导数为 $0$）。则函数 $z=f(x)\ln f(y)$ 在点 $(0,0)$ 处取得极小值的一个充分条件是（\quad）
  \begin{enumerate}
    \item[A.] $f(0)>1,\ f''(0)>0$；
    \item[B.] $f(0)>1,\ f''(0)<0$；
    \item[C.] $f(0)<1,\ f''(0)>0$；
    \item[D.] $f(0)<1,\ f''(0)<0$。
  \end{enumerate}
  \item 函数 $f(x,y)=1+x+y$ 在区域 $\{(x,y)\mid x^2+y^2\le 1\}$ 上的最大值与最小值之积是（\quad）
  \begin{enumerate}
    \item[A.] $-1$；
    \item[B.] $1$；
    \item[C.] $3-\sqrt2$；
    \item[D.] $1+\sqrt2$。
  \end{enumerate}
  \item 设 $\Gamma$ 是以 $(1,1),(-1,1),(-1,-1),(1,-1)$ 为顶点的正方形边界，则第一类曲线积分
  \[
    \oint_\Gamma \frac{x+y+1}{|x|+|y|}\,ds
  \]
  等于（\quad）
  \begin{enumerate}
    \item[A.] $0$；
    \item[B.] $8$；
    \item[C.] $4\ln 2$；
    \item[D.] $8\ln 2$。
  \end{enumerate}
  \item 设 $a$ 为常数，则级数
  \[
    \sum_{n=1}^{\infty}\left(\frac{\sin(an)}{n^2}-\frac1{\sqrt n}\right)
  \]
  （\quad）
  \begin{enumerate}
    \item[A.] 发散；
    \item[B.] 绝对收敛；
    \item[C.] 条件收敛；
    \item[D.] 敛散性与 $a$ 的取值相关。
  \end{enumerate}
\end{enumerate}

\paragraph*{三、计算题：共 3 题，每题 10 分，共 30 分。}
\begin{enumerate}
  \item 设
  \[
    g(x,y)=f\left(xy,\frac12(x^2-y^2)\right),
  \]
  其中 $f(u,v)$ 具有连续的二阶偏导数，且满足
  \[
    \frac{\partial^2 f}{\partial u^2}+\frac{\partial^2 f}{\partial v^2}=1.
  \]
  计算二阶偏导数 $\displaystyle \frac{\partial^2 g}{\partial x^2}+\frac{\partial^2 g}{\partial y^2}$。
  \item 计算累次积分
  \[
    \int_0^1 dx\int_0^x dy\int_0^y \frac{\sin z}{(1-z)^2}\,dz.
  \]
  \item 计算二重积分
  \[
    \iint_D \sqrt{x^2+y^2}\,dx\,dy,
  \]
  其中 $D$ 是由 $x=\sqrt{2y-y^2}$ 与 $y=x$ 所围成的闭区域。
\end{enumerate}

\paragraph*{四、综合题：共 3 题，每题 10 分，共 30 分。}
\begin{enumerate}
  \item 设 $\Sigma$ 是曲面 $z=1-x^2-y^2\ (z\ge 0)$ 的上侧，计算第二型曲面积分
  \[
    \iint_\Sigma 2x^3\,dy\,dz+2y^3\,dz\,dx+3(z^2-1)\,dx\,dy.
  \]
  \item 计算曲线积分
  \[
    \oint_L yz\,dx+3zx\,dy-xy\,dz,
  \]
  其中 $L$ 是曲线
  \[
    \begin{cases}
      x^2+y^2=2y,\\
      y-z=1,
    \end{cases}
  \]
  其方向从 $z$ 轴正向往 $z$ 轴负向看去为逆时针方向。
  \item 将函数
  \[
    f(x)=\frac{x}{x^2-5x+6}
  \]
  展开成 $x-1$ 的幂级数，并指出其收敛域。
\end{enumerate}

\paragraph*{五、证明题：共 1 题，每题 10 分，共 10 分。}
证明函数项级数
\[
  \sum_{n=1}^{\infty}\arctan\frac{2x}{x^2+n^3}
\]
在区间 $(-\infty,+\infty)$ 上一致收敛。

\paragraph*{六、解答题：共 1 题，每题 10 分，共 10 分。}
限制点 $(x,y)$ 在圆周 $(x-1)^2+y^2=1$ 上变化时，求函数 $f(x,y)=xy$ 的最小值与最大值。

\subsection{答案与解析}

\paragraph*{一、填空题}
\begin{enumerate}
  \item 特征方程 $r^2+3r+2=0$，得 $r=-1,-2$，故通解为
  \[
    y=C_1e^{-x}+C_2e^{-2x}.
  \]
  \item $\nabla u=(2y,2x,-2z)$，在 $(2,-1,-1)$ 处为 $(-2,4,2)$，方向导数最大值为
  \[
    |\nabla u|=\sqrt{(-2)^2+4^2+2^2}=2\sqrt6.
  \]
  \item 令 $x=\cos t,\ y=\sin t,\ 0\le t\le \dfrac{\pi}{2}$，则
  \[
    \int_\Gamma (x+y)^2\,ds
    =\int_0^{\pi/2}(\cos t+\sin t)^2\,dt
    =\frac{\pi}{2}+1.
  \]
  \item
  \[
    \sum_{n=0}^{\infty}(-1)^n\frac{2n+3}{(2n+1)!}
    =
    \sum_{n=0}^{\infty}\frac{(-1)^n}{(2n)!}
    +2\sum_{n=0}^{\infty}\frac{(-1)^n}{(2n+1)!}
    =\cos 1+2\sin 1.
  \]
  \item 在 $x=-\pi$ 处，周期延拓函数的右极限为 $\pi^2$，左极限为 $0$，故 Fourier 级数收敛于
  \[
    \frac{\pi^2+0}{2}=\frac{\pi^2}{2}.
  \]
\end{enumerate}

\paragraph*{二、选择题}
\begin{enumerate}
  \item 选 A。A 为二阶线性非齐次常微分方程；B 为一阶方程；C 含 $(y'')^2$，D 含 $\ln y$，均非二阶线性方程。
  \item 选 A。由 $f'(0)=0$，点 $(0,0)$ 为驻点。Hessian 矩阵对角元为 $f''(0)\ln f(0)$ 与 $f''(0)$，交叉项为 $0$；当 $f(0)>1,\ f''(0)>0$ 时正定，故取得极小值。
  \item 选 A。在线性函数 $1+x+y$ 在单位圆盘上，最大值为 $1+\sqrt2$，最小值为 $1-\sqrt2$，二者乘积为 $-1$。
  \item 选 D。沿正方形四条边分段计算，位于上边与右边的积分各为 $4\ln2$，下边与左边的积分和为 $0$，故总积分为 $8\ln2$。
  \item 选 A。$\sum \sin(an)/n^2$ 绝对收敛，而 $\sum 1/\sqrt n$ 发散，故原级数发散。
\end{enumerate}

\paragraph*{三、计算题}
\begin{enumerate}
  \item 记 $u=xy,\ v=\dfrac12(x^2-y^2)$，并用 $f_1,f_2$ 表示 $f$ 对第一、第二个变量的偏导数。则
  \[
    g_x=yf_1+xf_2,\qquad g_y=xf_1-yf_2.
  \]
  继续求导得
  \[
    g_{xx}=y^2f_{11}+2xyf_{12}+x^2f_{22}+f_2,
  \]
  \[
    g_{yy}=x^2f_{11}-2xyf_{12}+y^2f_{22}-f_2.
  \]
  因此
  \[
    g_{xx}+g_{yy}=(x^2+y^2)(f_{11}+f_{22})=x^2+y^2.
  \]
  \item 积分区域为 $0\le z\le y\le x\le 1$。交换积分次序得
  \[
    \int_0^1\frac{\sin z}{(1-z)^2}
    \left(\int_z^1dy\int_y^1dx\right)dz
    =
    \int_0^1\frac{\sin z}{(1-z)^2}\cdot\frac{(1-z)^2}{2}\,dz.
  \]
  故原式为
  \[
    \frac12\int_0^1\sin z\,dz=\frac{1-\cos 1}{2}.
  \]
  \item 曲线 $x=\sqrt{2y-y^2}$ 等价于 $x^2+(y-1)^2=1$，在极坐标下边界为 $r=2\sin\theta$；直线 $y=x$ 为 $\theta=\dfrac{\pi}{4}$。故
  \[
    D:\quad 0\le \theta\le \frac{\pi}{4},\quad 0\le r\le 2\sin\theta.
  \]
  于是
  \[
    \iint_D\sqrt{x^2+y^2}\,dx\,dy
    =
    \int_0^{\pi/4}\int_0^{2\sin\theta}r^2\,dr\,d\theta
    =
    \frac{16-10\sqrt2}{9}.
  \]
\end{enumerate}

\paragraph*{四、综合题}
\begin{enumerate}
  \item 用平面 $\Sigma_1:z=0,\ x^2+y^2\le 1$ 补成闭曲面，取外侧方向。由高斯公式，
  \[
    \iint_{\Sigma+\Sigma_1}2x^3\,dy\,dz+2y^3\,dz\,dx+3(z^2-1)\,dx\,dy
    =
    \iiint_\Omega (6x^2+6y^2+6z)\,dV=2\pi.
  \]
  在补面 $\Sigma_1$ 上取下侧，积分为
  \[
    I_{\Sigma_1}=-\iint_{x^2+y^2\le 1}3(z^2-1)\,dx\,dy=3\pi.
  \]
  因此所求曲面积分为
  \[
    2\pi-3\pi=-\pi.
  \]
  \item 按题设方向可取参数
  \[
    x=\cos t,\qquad y=1+\sin t,\qquad z=\sin t,\qquad 0\le t\le 2\pi.
  \]
  代入得
  \[
    \oint_L yz\,dx+3zx\,dy-xy\,dz
    =
    \int_0^{2\pi}\left[-(1+\sin t)\sin^2t+3\sin t\cos^2t-(1+\sin t)\cos^2t\right]dt.
  \]
  利用一个周期内奇函数项积分为 $0$，得
  \[
    -\int_0^{2\pi}\sin^2t\,dt-\int_0^{2\pi}\cos^2t\,dt=-2\pi.
  \]
  \item 分解为
  \[
    \frac{x}{x^2-5x+6}=\frac{3}{x-3}-\frac{2}{x-2}.
  \]
  令 $t=x-1$，则
  \[
    \frac{3}{x-3}=-\frac32\cdot\frac1{1-t/2}
    =-\frac32\sum_{n=0}^{\infty}\left(\frac{t}{2}\right)^n,\quad |t|<2,
  \]
  \[
    -\frac{2}{x-2}=\frac2{1-t}
    =2\sum_{n=0}^{\infty}t^n,\quad |t|<1.
  \]
  因此
  \[
    f(x)=\sum_{n=0}^{\infty}\left(2-\frac{3}{2^{n+1}}\right)(x-1)^n,
    \qquad |x-1|<1.
  \]
  收敛域为 $(0,2)$。
\end{enumerate}

\paragraph*{五、证明题}
对任意 $x\in\mathbb R$，由 $|\arctan u|\le |u|$ 以及
\[
  x^2+n^3\ge 2|x|n^{3/2}
\]
可得
\[
  \left|\arctan\frac{2x}{x^2+n^3}\right|
  \le
  \frac{2|x|}{x^2+n^3}
  \le
  \frac1{n^{3/2}}.
\]
而级数 $\sum_{n=1}^{\infty}n^{-3/2}$ 收敛，故由 Weierstrass 判别法知
\[
  \sum_{n=1}^{\infty}\arctan\frac{2x}{x^2+n^3}
\]
在 $(-\infty,+\infty)$ 上一致收敛。

\paragraph*{六、解答题}
令约束函数
\[
  \varphi(x,y)=(x-1)^2+y^2-1=0.
\]
对 $F(x,y,\lambda)=xy+\lambda\varphi(x,y)$，有
\[
  y+2\lambda(x-1)=0,\qquad x+2\lambda y=0.
\]
由约束得 $x^2+y^2=2x$。当 $\lambda=0$ 时得点 $(0,0)$，函数值为 $0$，不是最值点。设 $y\ne0$，由第二个方程得
\[
  \lambda=-\frac{x}{2y}.
\]
代入第一个方程，得
\[
  y^2=x(x-1).
\]
再与 $y^2=2x-x^2$ 联立，得
\[
  x(x-1)=2x-x^2,\qquad x=\frac32,\qquad y=\pm\frac{\sqrt3}{2}.
\]
因此
\[
  f\left(\frac32,\frac{\sqrt3}{2}\right)=\frac{3\sqrt3}{4}
  \]
为最大值，
\[
  f\left(\frac32,-\frac{\sqrt3}{2}\right)=-\frac{3\sqrt3}{4}
  \]
为最小值。
